Емкость проводников (холостой ход) $C_0 = $ \rule{2cm}{0.15mm} пФ.

\begin{table}[H]
    \centering
    \caption{Температурные зависимости емкости и $\tan \delta$ исследуемых образцов}
    \label{tab:measurements}
    \begin{tblr}[expand=\directlua]{
        width = \linewidth,
        colspec = { Q[c,m] *{10}{X[c,m]} },
        hlines, vlines,
        rows = {m, font=\small},
        rowhead = 2,
        stretch = 1.1
    }
        \SetCell[r=2]{c} $t$, $^\circ\text{C}$ & \SetCell[c=2]{c} {1. Неорг.\\ стекло} & & \SetCell[c=2]{c} {2. Слюда} & & \SetCell[c=2]{c} {3. Тиконд} & & \SetCell[c=2]{c} {4. Полипр-н} & & \SetCell[c=2]{c} {5. Сегнето-\\керамика} & \\
        & $C_1$, пФ & $\tan \delta_1$ & $C_2$, пФ & $\tan \delta_2$ & $C_3$, пФ & $\tan \delta_3$ & $C_4$, пФ & $\tan \delta_4$ & $C_5$, пФ & $\tan \delta_5$ \\
        & & & & & & & & & & \\
        & & & & & & & & & & \\
        & & & & & & & & & & \\
        & & & & & & & & & & \\
        & & & & & & & & & & \\
        & & & & & & & & & & \\
        & & & & & & & & & & \\
        & & & & & & & & & & \\
        & & & & & & & & & & \\
        & & & & & & & & & & \\
        & & & & & & & & & & \\
        & & & & & & & & & & \\
        & & & & & & & & & & \\
        & & & & & & & & & & \\
    \end{tblr}
\end{table}